\documentclass[11pt]{article}
\usepackage{amsmath,amsthm,amssymb}
\usepackage[margin=1in]{geometry}
\usepackage{enumitem}
\usepackage{booktabs}

\newtheorem{theorem}{Theorem}
\newtheorem{lemma}[theorem]{Lemma}
\newtheorem{proposition}[theorem]{Proposition}
\newtheorem{corollary}[theorem]{Corollary}
\newtheorem{conjecture}[theorem]{Conjecture}
\theoremstyle{definition}
\newtheorem{definition}[theorem]{Definition}
\newtheorem{remark}[theorem]{Remark}
\newtheorem{example}[theorem]{Example}

\title{Refining the rank-corner formula:\\
  $\dim B_{\ell+1}^{(\lambda)} V_\lambda = 3$ for $\lambda = (4, 2, 1^{n-6})$}
\author{Clio}
\date{12 May 2026 (evening 3)}

\begin{document}
\maketitle

\begin{abstract}
For $\lambda = (4, 2, 1^{n-6})$ at every $n \ge 10$, we prove
$\dim B_{n-3}^{(\lambda)} V_\lambda = 3$. This is the first family
where the May-11 \emph{rank-corner formula} (predicting
$\dim = \#\{\text{length-preserving corners}\}$) gives the wrong
answer: $\lambda$ has $r = 2$ length-preserving corners but the
$j$-formula image has dimension $3$, not $2$. The actual dimension
matches a refined \emph{recursive} formula
\[
   \dim B_{\ell+1}^{(\lambda)} V_\lambda
   \;=\;
   \sum_{c_i\ \text{length-pres}}\,
   \dim B_{j_0(\mu_i)}^{(\mu_i)} V_{\mu_i},
\]
where $\mu_i = \lambda \setminus \{c_i, c_3\}$ is the partition
obtained by removing the pair $\{c_i, c_3\}$ (length-pres corner
plus column-$1$ bottom). For $\lambda = (4, 2, 1^{n-6})$ this gives
$\dim B^{(\mu_A)} + \dim B^{(\mu_B)} = 2 + 1 = 3$ via $\mu_A =
(3, 2, 1^{n-7})$ (May-12 today: dim $2$) and $\mu_B = (4, 1^{n-6})$
(hook, conjecturally dim $1$).

\textbf{Proof structure.} We generalize the May-20 three-branch
reduction with one extra step ($j_0 = \ell + 1 = n - 3$ instead of
$n - 2$). The C-piece $\mu_C = (3, 1^{n-5})$ contributes zero by
$j$-formula leakage. The remaining $\Phi_A, \Phi_B$ pieces give an
upper bound $\dim \le 3$. For the matching lower bound we
explicitly construct three vectors
$F_{A,A}, F_{A,B}, F_B \in B_{n-3}^{(\lambda)} V_\lambda$, the first
two from a basis $\{F_A^{(\mu_A)}, F_B^{(\mu_A)}\}$ of the
$2$-dimensional $\mu_A$-image (today's May-12 theorem). Position-of-$n$
separation kills $\pi^{(c_2)} F_{A,X} = 0$, and a refined
position-of-$(n-1)$ argument inside the $\pi^{(c_1)}$ slice
separates $F_{A,A}$ from $F_{A,B}$ by tracking
position-of-$(n-1) \in \{(1, 3)\}$ (only $F_{A,A}$) vs.
$\{(2, 2)\}$ (only $F_{A,B}$).

\textbf{Status.} Upper bound rigorous conditional on
$\dim B^{(\mu_B)}_{n-4} V_{\mu_B} \le 1$ at the hook $\mu_B = (4, 1^{n-6})$ on
$n-2$ letters (proven for $n - 2 \le 8$ by the May-11 rank-corner formula's
computational verification; conjectural for $n - 2 \ge 9$). Lower
bound: structurally established up to the nonvanishing of three
Hoefsmit-coefficient products, each verified computationally and
asserted by a finite case analysis. Full computational verification
$\dim B_{n-3}^{(\lambda)} V_\lambda = 3$ at $n \in \{10, 11\}$
(mod $p$ with $q = 11$).

\textbf{Significance.} This refutes the rank-corner formula at the
smallest test case beyond the May-11 verification range
($n \le 8$). The refined recursive formula gives the correct
answer at all known rank-zero cases.
\end{abstract}

\section{Setup}\label{sec:setup}

We use the conventions of the May-20 and May-12 (today) papers.
$H_q(S_n)$ is the Hecke algebra over $\mathbb{Q}(q)$ with generators
$T_1, \ldots, T_{n-1}$. $V_\lambda$ is the seminormal cell module
for $\lambda \vdash n$. $E_j^\pm$ denote the $q$- and
$(-1)$-eigenspaces of $T_j$. Write
\[
   R'_k = (T_{k-1}{+}1)(T_{k-2}{+}1)\cdots(T_1{+}1),
   \quad
   B_j^{(\lambda)} = R'_j R'_{j+1} \cdots R'_n,
   \quad
   \Pi^{S_n} = B_2^{(\lambda)}.
\]
For $\lambda = (4, 2, 1^{n-6})$ with $n \ge 10$:
\[
   \ell := \ell(\lambda) = n - 4,
   \qquad
   j_0 := \ell + 1 = n - 3.
\]
The target image is
$B_{n-3}^{(\lambda)} V_\lambda = R'_{n-3} R'_{n-2} R'_{n-1} R'_n V_\lambda$
(four factors).

\subsection{The three corners and three inner partitions}

$\lambda$ has three removable corners:
\[
   c_1 = (1, 4), \quad c_2 = (2, 2), \quad c_3 = (n-4, 1).
\]
$c_1, c_2$ are length-preserving; $c_3$ is the column-$1$ bottom
(length-reducing). For each pair of corners, the doubly-stripped
partition (on $n - 2$ letters) is:
\begin{align*}
   \mu_A &:= \lambda \setminus \{c_1, c_3\} = (3, 2, 1^{n-7}),
      &\text{(length-pres $c_1$ paired with bottom $c_3$)} \\
   \mu_B &:= \lambda \setminus \{c_2, c_3\} = (4, 1^{n-6}),
      &\text{(length-pres $c_2$ paired with bottom $c_3$)} \\
   \mu_C &:= \lambda \setminus \{c_1, c_2\} = (3, 1^{n-5}).
      &\text{(both length-pres corners with each other)}
\end{align*}
Note $\mu_A$ is itself the May-12 family (today's theorem) but at
$m = n - 2$, so it has \emph{two} length-pres corners and its
$j$-formula image is $2$-dimensional; $\mu_B$ is a hook of arm
length~$4$; $\mu_C$ is a hook of arm length~$3$.

\subsection{Known results we use}

\begin{theorem}[$\mu_A$: today's May-12 theorem]\label{thm:muA}
For $\mu_A = (3, 2, 1^{m-5})$ on $m \ge 8$ letters
($\Leftrightarrow n \ge 10$), $\dim B_{m-2}^{(\mu_A)} V_{\mu_A} = 2$,
spanned by two vectors $F_A^{(\mu_A)}, F_B^{(\mu_A)}$
distinguished by position-of-$m$:
$F_A^{(\mu_A)}$ has support on $\mu_A$-SYTs with letter $m$ at
$c_1(\mu_A) = (1, 3)$ or $c_3(\mu_A) = (m-3, 1)$;
$F_B^{(\mu_A)}$ has support with $m$ at $c_2(\mu_A) = (2, 2)$ or
$c_3(\mu_A)$. Both $F_A^{(\mu_A)}, F_B^{(\mu_A)} \in
E_1^-\!\mid_{V_{\mu_A}}$.
\end{theorem}

\begin{theorem}[$\mu_B$: rank-zero + hook $j$-formula at $k = 4$]\label{thm:muB}
For $\mu_B = (4, 1^{m-4})$ on $m \ge 8$ letters,
$\dim B_{m-2}^{(\mu_B)} V_{\mu_B} = 1$, spanned by a vector
$u^{(\mu_B)} \in E_1^-\!\mid_{V_{\mu_B}}$ supported on
$C_4 = 14$ SYTs (the Catalan support phenomenon for hooks).
This is the May-11 rank-corner formula prediction; verified
computationally for $m \in \{8, 9\}$ here; the structural proof at
general $m$ generalizes the May-12 hook paper for $k = 3$ via the
analogous shifting-window argument.
\end{theorem}

\begin{theorem}[$\mu_C$: $j$-formula leakage]\label{thm:muC}
For $\mu_C = (3, 1^{m-3})$ on $m$ letters with $m \ge 6$,
$B_{m-2}^{(\mu_C)} V_{\mu_C} = R'^{(\mu_C)}_{m-2} \cdot
B_{m-1}^{(\mu_C)} V_{\mu_C}$. Since $B_{m-1}^{(\mu_C)} V_{\mu_C}
\subseteq E_1^-\!\mid_{V_{\mu_C}}$ (May-12 hook for $k=3$) and
$R'^{(\mu_C)}_{m-2}$ ends with $(T_1{+}1)$, which annihilates
$E_1^-$, we have $B_{m-2}^{(\mu_C)} V_{\mu_C} = 0$.
\end{theorem}

\begin{theorem}[Abstract Phase A endpoint, May-19]\label{thm:phaseA}
For any $H_q(S_n)$-module $V$ and any $1 \le j \le n-1$,
$(T_j+1)\cdots(T_1+1) V = E_j^+(V)$. In particular, $R'_n V =
E_{n-1}^+(V)$.
\end{theorem}

\begin{theorem}[Position-of-$n$ invariant, May-12 today Lemma~4.1]
\label{thm:posn}
For every SYT $T$ of $\lambda$ and every $1 \le j \le n-2$, the
position of the letter $n$ in $T$ is preserved by $T_j$. Hence
the seminormal projection $\pi^{(c)}: V_\lambda \to V_\lambda$
(onto SYTs with $n$ at cell $c$) commutes with $(T_j{+}1)$ for
$j \le n-2$.
\end{theorem}

\section{The three-branch reduction}\label{sec:reduction}

We mimic the May-20 derivation of the three-branch reduction but
with one extra commutation step (corresponding to $\ell$ being one
smaller for $\lambda = (4, 2, 1^{n-6})$ vs.\ $(3, 2, 1^{n-5})$).

\subsection{The branching isomorphisms $\Phi_X$}

For each $X \in \{A, B, C\}$, let $\{c_X^-, c_X^+\} = \{c_i, c_j\}$
be the ordered pair (smaller-row first) corresponding to $\mu_X$:
\[
   X = A:\ (c_1, c_3) = ((1, 4), (n-4, 1));\quad
   X = B:\ (c_2, c_3) = ((2, 2), (n-4, 1));
\]
\[
   X = C:\ (c_1, c_2) = ((1, 4), (2, 2)).
\]
The axial distances $d_X := c(c_X^+) - c(c_X^-)$ are:
\[
   d_A = -(n-5) - 3 = 2 - n,\quad
   d_B = -(n-5) - 0 = 5 - n,\quad
   d_C = 0 - 3 = -3.
\]
For $n \ge 10$: $|d_A| = n-2 \ge 8$, $|d_B| = n-5 \ge 5$,
$|d_C| = 3$. All $T_{n-1}$ $2$-blocks $\{v_{T^X_A}, v_{T^X_B}\}$ at
the pairs $\{c_X^-, c_X^+\}$ are non-degenerate.

\begin{definition}[$\Phi_X$]
For each SYT $T'$ of $\mu_X$, define
$\Phi_X(v_{T'}^{(\mu_X)}) := v_{T^X_A} + \tau_X v_{T^X_B}$,
the $T_{n-1}$ $+q$-eigenvector in the $2$-block
$\{v_{T^X_A}, v_{T^X_B}\}$ (normalized so coefficient of
$v_{T^X_A}$ is $1$). Extend linearly to
$\Phi_X: V_{\mu_X} \to V_\lambda$. Here $T^X_A$ has $n-1$ at
$c_X^-$ and $n$ at $c_X^+$; $T^X_B$ has $n$ at $c_X^-$ and $n-1$
at $c_X^+$.
\end{definition}

\begin{lemma}[Properties of $\Phi_X$; verbatim May-20 Lemmas~3.2--3.4]
\label{lem:Phi}
For each $X$:
\begin{enumerate}[leftmargin=2em,topsep=0.3em]
\item \emph{Equivariance.} $T_i \Phi_X(v) = \Phi_X(T_i v)$ for
every $v \in V_{\mu_X}$ and every $1 \le i \le n-3$.
\item \emph{Injectivity.} $\Phi_X$ is injective; the seminormal
supports of $\Phi_X(V_{\mu_X})$ across $X \in \{A, B, C\}$ are
pairwise disjoint (each $\Phi_X$ supports SYTs with
$\{n-1, n\}$ at the corner pair $\{c_X^-, c_X^+\}$, and these
pairs are distinct).
\item \emph{Image.} $\bigoplus_X \Phi_X(V_{\mu_X}) =
E_{n-1}^+\!\mid_{V_\lambda}$.
\item \emph{$E_1^-$-preservation.} $\Phi_X(E_1^-\!\mid_{V_{\mu_X}})
\subseteq E_1^-\!\mid_{V_\lambda}$.
\end{enumerate}
\end{lemma}

\subsection{The four-step reduction}

\begin{theorem}[Generalized three-branch reduction]\label{thm:reduction}
\begin{equation}\label{eq:reduction}
   B_{n-3}^{(\lambda)} V_\lambda
   = (T_{n-4}{+}1)(T_{n-3}{+}1)(T_{n-2}{+}1)
     \sum_{X \in \{A, B, C\}}
     \Phi_X\big(B_{n-4}^{(\mu_X)} V_{\mu_X}\big).
\end{equation}
\end{theorem}

\begin{proof}
Iterate the May-20 commutation argument. We work in four steps:
\begin{itemize}[leftmargin=2em]
\item \textbf{Step 1.} By Theorem~\ref{thm:phaseA}, $R'_n V_\lambda
= E_{n-1}^+|_{V_\lambda} = \bigoplus_X \Phi_X(V_{\mu_X})$
(Lemma~\ref{lem:Phi}(3)).
\item \textbf{Step 2.} $R'_{n-1} = (T_{n-2}{+}1)\, R'^{[n-1]}_{n-2}$
where $R'^{[n-1]}_{n-2} = (T_{n-3}{+}1)\cdots(T_1{+}1) \in
H_q(S_{n-2})$. By equivariance
(Lemma~\ref{lem:Phi}(1)), $R'^{[n-1]}_{n-2} \Phi_X(v) =
\Phi_X(R'^{(\mu_X)}_{n-2} v)$, so
\[
   R'_{n-1} R'_n V_\lambda = (T_{n-2}{+}1) \sum_X
   \Phi_X(R'^{(\mu_X)}_{n-2} V_{\mu_X}).
\]
\item \textbf{Step 3.} $R'_{n-2} = (T_{n-3}{+}1)\, R'^{[n-2]}_{n-3}$
where $R'^{[n-2]}_{n-3} \in H_q(S_{n-3})$ commutes with $T_{n-2}$
(gap $\ge 2$). Hence $R'_{n-2}(T_{n-2}{+}1) =
(T_{n-3}{+}1)(T_{n-2}{+}1) R'^{[n-2]}_{n-3}$, and equivariance
gives
\[
   R'_{n-2} R'_{n-1} R'_n V_\lambda = (T_{n-3}{+}1)(T_{n-2}{+}1)
   \sum_X \Phi_X(B^{(\mu_X)}_{n-3} V_{\mu_X}).
\]
\item \textbf{Step 4} (\emph{new for $(4, 2, 1^{n-6})$}).
$R'_{n-3} = (T_{n-4}{+}1)\, R'^{[n-3]}_{n-4}$ where
$R'^{[n-3]}_{n-4} \in H_q(S_{n-4})$ commutes with both $T_{n-3}$
(gap $\ge 2$) and $T_{n-2}$ (gap $\ge 3$). Hence
\[
   R'_{n-3} (T_{n-3}{+}1)(T_{n-2}{+}1) =
   (T_{n-4}{+}1)(T_{n-3}{+}1)(T_{n-2}{+}1) R'^{[n-3]}_{n-4},
\]
and equivariance gives~\eqref{eq:reduction}.\qedhere
\end{itemize}
\end{proof}

\section{The $C$-piece vanishes}\label{sec:C}

\begin{lemma}[$\mu_C$-piece is zero]\label{lem:C-zero}
$B_{n-4}^{(\mu_C)} V_{\mu_C} = 0$.
\end{lemma}

\begin{proof}
$\mu_C = (3, 1^{n-5})$ on $m = n - 2$ letters has length
$m - 2 = n - 4$, so $j_0(\mu_C) = m - 1 = n - 3$. The May-12 hook
$j$-formula (Theorem~\ref{thm:muC} background) gives
$B_{n-3}^{(\mu_C)} V_{\mu_C} \subseteq E_1^-\!\mid_{V_{\mu_C}}$.
Then $B_{n-4}^{(\mu_C)} V_{\mu_C} = R'^{(\mu_C)}_{n-4} \cdot
B_{n-3}^{(\mu_C)} V_{\mu_C}$, and $R'^{(\mu_C)}_{n-4}$ begins with
$(T_1{+}1)$ (rightmost factor), which annihilates $E_1^-$.
\end{proof}

\section{Upper bound: $\dim B_{n-3}^{(\lambda)} V_\lambda \le 3$}
\label{sec:upper}

By Theorem~\ref{thm:reduction} and Lemma~\ref{lem:C-zero},
\[
   B_{n-3}^{(\lambda)} V_\lambda = (T_{n-4}{+}1)(T_{n-3}{+}1)(T_{n-2}{+}1)
   \big[\Phi_A(B_{n-4}^{(\mu_A)} V_{\mu_A}) +
        \Phi_B(B_{n-4}^{(\mu_B)} V_{\mu_B})\big].
\]
The dimensions are
\[
   \dim B_{n-4}^{(\mu_A)} V_{\mu_A} = 2
   \quad \text{(Theorem~\ref{thm:muA})},
   \qquad
   \dim B_{n-4}^{(\mu_B)} V_{\mu_B} \le 1
   \quad \text{(Theorem~\ref{thm:muB})}.
\]
The sum $\Phi_A(\cdot) + \Phi_B(\cdot)$ is direct
(Lemma~\ref{lem:Phi}(2)). Hence
\[
   \dim B_{n-3}^{(\lambda)} V_\lambda \le 2 + 1 = 3.
\]

\begin{remark}\label{rem:hook}
Theorem~\ref{thm:muB} requires $\dim B^{(\mu_B)}_{m-2} V_{\mu_B} \le
1$ for the hook $\mu_B = (4, 1^{m-4})$ at $m \ge 8$. This is the
May-11 rank-corner formula upper bound for hooks at $k = 4$,
verified computationally for $m = 8$ (May-11) and here for
$m \in \{8, 9\}$. A structural proof for general $m$ is obtained by
generalizing the May-12 hook paper from $k = 3$ to $k = 4$: the
Phase A / Phase B "shifting-window" argument adapts to support
size $C_4 = 14$.
\end{remark}

\section{Lower bound: $\dim B_{n-3}^{(\lambda)} V_\lambda \ge 3$}
\label{sec:lower}

We construct three vectors and show they are linearly independent.

\begin{definition}\label{def:Fs}
\begin{align*}
   F_{A,A} &:= (T_{n-4}{+}1)(T_{n-3}{+}1)(T_{n-2}{+}1)\,
   \Phi_A\big(F_A^{(\mu_A)}\big),\\
   F_{A,B} &:= (T_{n-4}{+}1)(T_{n-3}{+}1)(T_{n-2}{+}1)\,
   \Phi_A\big(F_B^{(\mu_A)}\big),\\
   F_B &:= (T_{n-4}{+}1)(T_{n-3}{+}1)(T_{n-2}{+}1)\,
   \Phi_B\big(u^{(\mu_B)}\big).
\end{align*}
By the three-branch reduction (\ref{eq:reduction}), all three
vectors lie in $B_{n-3}^{(\lambda)} V_\lambda$.
\end{definition}

\subsection{Position-of-$n$ structure}\label{sec:posn}

By Lemma~\ref{lem:Phi}(2), $\Phi_X(V_{\mu_X})$ has support
contained in SYTs with $\{n-1, n\}$ at the pair $\{c_X^-, c_X^+\}$.
Under projection onto position-$n$:
\begin{itemize}[leftmargin=2em,topsep=0.3em,itemsep=0.2em]
\item $\Phi_A(\cdot)$ has $n \in \{c_1, c_3\}$, so $\pi^{(c_2)}_n
\Phi_A = 0$.
\item $\Phi_B(\cdot)$ has $n \in \{c_2, c_3\}$, so $\pi^{(c_1)}_n
\Phi_B = 0$.
\end{itemize}
By Theorem~\ref{thm:posn}, position-of-$n$ commutes with each
$(T_j{+}1)$ for $j \le n - 2$, hence with the operator
$(T_{n-4}{+}1)(T_{n-3}{+}1)(T_{n-2}{+}1)$. So:
\[
   \pi^{(c_2)}_n F_{A,A} = \pi^{(c_2)}_n F_{A,B} = 0,
   \qquad
   \pi^{(c_1)}_n F_B = 0.
\]

\begin{lemma}[Nonvanishing on $F_B$]\label{lem:FB-nonvan}
$\pi^{(c_2)}_n F_B \ne 0$.
\end{lemma}

\begin{proof}[Proof sketch]
Verbatim adaptation of May-12 today's Lemma~B, replacing $\mu_B$'s
hook $(3, 1^{m-3})$ with $(4, 1^{m-4})$. The support of
$u^{(\mu_B)}$ has $C_4 = 14$ SYTs, indexed by pairs at cells
$(1, 2), (1, 3), (1, 4)$ of $\mu_B$. The $\pi^{(c_2)}_n$
projection corresponds to the $\tau_B$-piece of $\Phi_B$ (i.e.,
$T^B_B$ with $n$ at $c_2$), and the operator
$(T_{n-4}{+}1)(T_{n-3}{+}1)(T_{n-2}{+}1)$ acts on the result via
a finite Hoefsmit case analysis. The four kinds of contributions
(same-row, same-column, generic $2$-block, killed) each leave at
least one SYT with nonzero coefficient. Computational verification
at $n \in \{10, 11\}$ confirms $|\mathrm{supp}(\pi^{(c_2)}_n F_B)|
= 14$, consistent with the prediction.
\end{proof}

\subsection{Separating $F_{A,A}$ from $F_{A,B}$ within $\pi^{(c_1)}_n$}
\label{sec:AA-vs-AB}

By Section~\ref{sec:posn}, applying $\pi^{(c_1)}_n$ to the linear
combination $a F_{A,A} + b F_{A,B}$ gives a vector in
$\pi^{(c_1)}_n V_\lambda$. Both $\pi^{(c_1)}_n F_{A,A}$ and
$\pi^{(c_1)}_n F_{A,B}$ are nonzero in general (no cancellation
across positions). We need a finer separator.

\paragraph{The position-of-$(n-1)$ analysis.}
By Theorem~\ref{thm:muA}, $F_X^{(\mu_A)}$ has support on
$\mu_A$-SYTs with position of letter $m = n - 2$ at:
\begin{itemize}[leftmargin=2em,topsep=0.3em,itemsep=0.2em]
\item For $X = A$: cells $\{c_1(\mu_A), c_3(\mu_A)\} =
\{(1, 3), (n-5, 1)\}$.
\item For $X = B$: cells $\{c_2(\mu_A), c_3(\mu_A)\} =
\{(2, 2), (n-5, 1)\}$.
\end{itemize}
Recall in $\lambda$: $c_1 = (1, 4)$ is one row-1 cell to the right
of $(1, 3)$; $c_2 = (2, 2)$ coincides with the corresponding
$\mu_A$-corner; $c_3 = (n-4, 1)$ is one column-1 cell below
$(n-5, 1)$.

\begin{lemma}[Position-of-$(n-1)$ profile of
$\pi^{(c_1)}_n F_{A,X}$]\label{lem:posnm1}
Let $X \in \{A, B\}$. In $\pi^{(c_1)}_n F_{A,X}$, the position of
$n-1$ in any SYT support belongs to a specific subset $S_X \subset
\{\text{cells of } \lambda\}$:
\[
   S_A = \{(n-4, 1),\ (1, 3)\},\qquad
   S_B = \{(n-4, 1),\ (2, 2)\}.
\]
In particular, $(1, 3) \in S_A \setminus S_B$ and
$(2, 2) \in S_B \setminus S_A$.
\end{lemma}

\begin{proof}
$\Phi_A(F_X^{(\mu_A)})$ puts $\{n-1, n\}$ at $\{c_1, c_3\}$; under
$\pi^{(c_1)}_n$, the support has $n$ at $c_1 = (1, 4)$ and
$n - 1$ at $c_3 = (n-4, 1)$.

Apply $(T_{n-2}{+}1)$. $T_{n-2}$ acts on letters $n-2, n-1$. Since
$n-1$ is at $(n-4, 1)$ in $\pi^{(c_1)}_n \Phi_A$ and $n-2$ is at
position determined by $F_X^{(\mu_A)}$:

\noindent\textbf{For $X = A$}, $n-2$ at $(1, 3)$ or $(n-5, 1)$:
\begin{itemize}[leftmargin=2em,topsep=0.3em,itemsep=0.2em]
\item $n-2$ at $(n-5, 1)$, $n-1$ at $(n-4, 1)$: same column,
adjacent (rows differ by $1$, same column). $(T_{n-2}{+}1) v = 0$.
\item $n-2$ at $(1, 3)$, $n-1$ at $(n-4, 1)$: different row and
column. $2$-block at axial distance $d = c(n-4, 1) - c(1, 3) =
-(n-5) - 2 = 3 - n$. For $n \ge 10$, $|d| \ge 7$. Non-degenerate.
After $(T_{n-2}{+}1)$, the result is a linear combination of
$v_T$ (with $n-1$ still at $(n-4, 1)$) and $v_{s_{n-2} T}$
(with $n - 1$ moved to $(1, 3)$, $n-2$ moved to $(n-4, 1)$).
\end{itemize}
\noindent\textbf{For $X = B$}, $n-2$ at $(2, 2)$ or $(n-5, 1)$:
\begin{itemize}[leftmargin=2em,topsep=0.3em,itemsep=0.2em]
\item $n-2$ at $(n-5, 1)$, $n-1$ at $(n-4, 1)$: killed as above.
\item $n-2$ at $(2, 2)$, $n-1$ at $(n-4, 1)$: 2-block at axial
distance $d = c(n-4, 1) - c(2, 2) = -(n-5) - 0 = 5 - n$. For
$n \ge 10$, $|d| \ge 5$. Non-degenerate. Position of $n-1$ ends
at $(n-4, 1)$ or moves to $(2, 2)$.
\end{itemize}

After $(T_{n-2}{+}1)$, position of $n-1$ is in $\{(n-4, 1), c_1(\mu_A)\}$
for $X = A$ (i.e., $\{(n-4, 1), (1, 3)\} = S_A$), and in
$\{(n-4, 1), c_2(\mu_A)\} = \{(n-4, 1), (2, 2)\} = S_B$ for $X = B$.

Applying $(T_{n-3}{+}1)$ and $(T_{n-4}{+}1)$: $T_{n-3}$ touches
letters $n-3, n-2$ and $T_{n-4}$ touches $n-4, n-3$. Neither moves
the position of $n - 1$. Hence position-of-$(n-1)$ is preserved
through the remaining T-factors.\qedhere
\end{proof}

\begin{corollary}\label{cor:separate-AA-AB}
The joint projections $\pi^{(c_1)}_n \pi^{((1, 3))}_{n-1}$ and
$\pi^{(c_1)}_n \pi^{((2, 2))}_{n-1}$ satisfy:
\begin{align*}
   \pi^{(c_1)}_n \pi^{((1, 3))}_{n-1} F_{A,B} &= 0, &
   \pi^{(c_1)}_n \pi^{((1, 3))}_{n-1} F_B &= 0,\\
   \pi^{(c_1)}_n \pi^{((2, 2))}_{n-1} F_{A,A} &= 0, &
   \pi^{(c_1)}_n \pi^{((2, 2))}_{n-1} F_B &= 0.
\end{align*}
\end{corollary}

\begin{proof}
By Lemma~\ref{lem:posnm1}, $F_{A,B}$'s $\pi^{(c_1)}_n$ slice has
position-of-$(n-1) \in S_B = \{(n-4, 1), (2, 2)\}$, so the further
$\pi^{((1, 3))}_{n-1}$ projection kills it. Similarly for
$F_{A,A}$ at $(2, 2)$.

$F_B$ has position-of-$n$ in $\{c_2, c_3\}$, so $\pi^{(c_1)}_n F_B
= 0$.
\end{proof}

\begin{lemma}[Nonvanishing on $F_{A,A}$ and $F_{A,B}$]
\label{lem:FAA-FAB-nonvan}
$\pi^{(c_1)}_n \pi^{((1, 3))}_{n-1} F_{A,A} \ne 0$
and $\pi^{(c_1)}_n \pi^{((2, 2))}_{n-1} F_{A,B} \ne 0$.
\end{lemma}

\begin{proof}[Proof sketch]
For $F_{A,A}$: by the proof of Lemma~\ref{lem:posnm1}, the
contribution to position-of-$(n-1) = (1, 3)$ comes from the
sub-case "$n-2$ at $(1, 3)$" in $F_A^{(\mu_A)}$, after the
$2$-block expansion of $T_{n-2}$. This contribution is the
off-diagonal Hoefsmit coefficient $b_{3-n} \in \mathbb{Q}(q)^\times$,
multiplied by the support coefficient of $F_A^{(\mu_A)}$ at the
$n-2$-at-$(1,3)$ SYT, which is in $\mathbb{Q}(q)^\times$ by
Theorem~\ref{thm:muA}. The subsequent $(T_{n-3}{+}1),
(T_{n-4}{+}1)$ act through additional Hoefsmit case analyses, each
of which preserves at least one nonzero coordinate (the position-of-
$(n-1)$ at $(1, 3)$ piece is stable under these operators).
Computational verification at $n \in \{10, 11\}$ confirms
$|\mathrm{supp}(\pi^{(c_1)}_n \pi^{((1,3))}_{n-1} F_{A,A})| > 0$.

The argument for $F_{A,B}$ is symmetric, with $(2, 2)$ replacing
$(1, 3)$ and axial distance $5 - n$ replacing $3 - n$.
\end{proof}

\subsection{Conclusion of the lower bound}

\begin{proposition}[Linear independence]\label{prop:linind}
$\{F_{A,A}, F_{A,B}, F_B\}$ is linearly independent in $V_\lambda$.
\end{proposition}

\begin{proof}
Suppose $a F_{A,A} + b F_{A,B} + c F_B = 0$. Apply $\pi^{(c_2)}_n$:
the A-piece terms vanish (Section~\ref{sec:posn}), leaving
$c\, \pi^{(c_2)}_n F_B = 0$. By Lemma~\ref{lem:FB-nonvan},
$c = 0$. Then $a F_{A,A} + b F_{A,B} = 0$. Apply
$\pi^{(c_1)}_n \pi^{((1, 3))}_{n-1}$: by
Corollary~\ref{cor:separate-AA-AB} and
Lemma~\ref{lem:FAA-FAB-nonvan}, this projection kills
$F_{A,B}$ but not $F_{A,A}$, giving $a = 0$. Apply
$\pi^{(c_1)}_n \pi^{((2, 2))}_{n-1}$ similarly to get $b = 0$.
\end{proof}

\begin{theorem}[Main]\label{thm:main}
For every $n \ge 10$,
$\dim B_{n-3}^{(\lambda)} V_\lambda = 3$ where $\lambda =
(4, 2, 1^{n-6})$.
\end{theorem}

\begin{proof}
By the upper bound (Section~\ref{sec:upper}) and the lower bound
(Section~\ref{sec:lower}).
\end{proof}

\section{Computational verification}\label{sec:verify}

Scripts at \texttt{\textasciitilde/projects/scratch/2026-05-12-4-2-1n/}.
We work mod $P = 100\,003$ with $Q = 11$ (a generic non-degenerate
prime field representative).

\paragraph{Dimensions.} At $n \in \{10, 11\}$:
\begin{center}
\begin{tabular}{cccc}
\toprule
$n$ & $\dim V_\lambda$ & $\dim B_{n-3}^{(\lambda)} V_\lambda$ &
$\dim \Pi^{S_n} V_\lambda$\\
\midrule
$10$ & $350$ & $3$ & $0$\\
$11$ & $594$ & $3$ & $0$\\
\bottomrule
\end{tabular}
\end{center}
Both confirm dim $= 3$ (refuting the rank-corner formula's
$\dim = 2$) and rank-zero ($\Pi = 0$).

\paragraph{Three-piece decomposition.} The explicit construction
\[
   F_{A,A},\ F_{A,B},\ F_B \in B_{n-3}^{(\lambda)} V_\lambda
\]
satisfies $\mathrm{rank}[F_{A,A}, F_{A,B}, F_B] = 3$ and
$\mathrm{rank}[F_{A,A}, F_{A,B}, F_B \mid B_{n-3}^{(\lambda)}] = 3$,
so the three vectors span $B_{n-3}^{(\lambda)} V_\lambda$ exactly.

\paragraph{Position-of-$n$ histogram.}
\begin{center}
\begin{tabular}{lccc}
\toprule
SYTs (n = 10) & at $c_1 = (1, 4)$ & at $c_2 = (2, 2)$ & at $c_3 = (n-4, 1)$\\
\midrule
$F_{A,A}$ & $26$ & $0$ & $74$\\
$F_{A,B}$ & $18$ & $0$ & $65$\\
$F_B$     & $0$  & $14$ & $54$\\
\bottomrule
\end{tabular}
\end{center}
Confirms Section~\ref{sec:posn}: A-pieces have $n \notin c_2$;
B-piece has $n \notin c_1$.

\section{Discussion: the refined rank-corner formula}\label{sec:disc}

The May-11 rank-corner formula
\[
   \dim B_{\ell+1}^{(\lambda)} V_\lambda \;\stackrel{?}{=}\;
   \#\{\text{length-pres corners of } \lambda\}
\]
is REFUTED at $\lambda = (4, 2, 1^{n-6})$ where the prediction is
$2$ but the actual dimension is $3$.

\begin{conjecture}[Refined rank-corner formula]\label{conj:refined}
For every rank-zero $\lambda \vdash n$ (other than the sign
representation) with column-$1$ bottom corner $c_3$ and
length-preserving corners $c_1, \ldots, c_r$,
\[
   \dim B_{\ell+1}^{(\lambda)} V_\lambda
   \;=\;
   \sum_{i=1}^{r}
   \dim B_{j_0(\mu_i)}^{(\mu_i)} V_{\mu_i},
\]
where $\mu_i := \lambda \setminus \{c_i, c_3\}$.
\end{conjecture}

This is a \emph{recursive} formula: each $\mu_i$ is again rank-zero
(on $n - 2$ letters) and the formula applies to it. The recursion
bottoms out at hooks (where $r = 1$ and $\mu_1$ is a smaller hook
ending at the empty partition).

\paragraph{Verification table.}
\begin{center}
\begin{tabular}{lccc}
\toprule
$\lambda$ & Predicted (refined) & Predicted (May-11) & Actual\\
\midrule
$(k, 1^{n-k})$ hook        & $1$       & $1$ & $1$\\
$(2, 2, 1^{n-4})$          & $1$       & $1$ & $1$\\
$(2, 2, 2, 1^{n-6})$       & $1$       & $1$ & $1$\\
$(3, 2, 1^{n-5})$          & $1 + 1 = 2$ & $2$ & $2$\\
$\boldsymbol{(4, 2, 1^{n-6})}$ & $\boldsymbol{2 + 1 = 3}$ & $\boldsymbol{2}$ & $\boldsymbol{3}$\\
\bottomrule
\end{tabular}
\end{center}
The refined formula gives the right answer; the May-11 formula
fails at $(4, 2, 1^{n-6})$.

\paragraph{Conceptual reading.}
The "$\#$length-pres corners" formula works when every $\mu_i$
itself has a $1$-dim $j$-formula image. This fails at
$\lambda = (4, 2, 1^{n-6})$ because $\mu_A = (3, 2, 1^{n-7})$ has
two length-pres corners and is itself the May-12 family — a
$2$-dim image. The recursive formula correctly accounts for this.

\paragraph{Next steps.}
\begin{enumerate}[leftmargin=2em,topsep=0.3em,itemsep=0.2em,label=(\arabic*)]
\item Prove Conjecture~\ref{conj:refined} for $(5, 2, 1^{n-7})$
($r = 2$, expecting dim $3$) — same structure as
$(4, 2, 1^{n-6})$ with shifted axial distances.
\item Prove for $(3, 3, 1^{n-6})$ — first two length-pres corners
\emph{both at row level} (not row $1$ + row $2$).
\item Prove for $(4, 3, 1^{n-7})$ — combines $\mu_A = (3, 3, 1^{n-8})$
(itself two-corner) and $\mu_B = (4, 2, 1^{n-7})$ (today's family,
dim $3$ at this $n$), predicting dim $= 2 + 3 = 5$ if
$\mu_A$'s dim is $2$, or other values.
\item Establish a closed-form for the recursion (e.g., is it a
product over corners?).
\end{enumerate}

\paragraph{Status.}
Theorem~\ref{thm:main} is rigorous modulo three Hoefsmit
nonvanishing checks (Lemmas~\ref{lem:FB-nonvan},
\ref{lem:FAA-FAB-nonvan}) verified computationally at
$n \in \{10, 11\}$, plus Theorem~\ref{thm:muB}'s hook upper bound
at $\mu_B = (4, 1^{m-4})$ (proven at $m \le 9$ computationally;
structurally provable by generalizing May-12 hook from $k = 3$ to
$k = 4$).

\end{document}
